\documentclass[11pt,a4paper]{article}

\usepackage[margin=1in]{geometry}
\usepackage{amsmath,amssymb,amsfonts}
\usepackage{mathtools}
\usepackage{booktabs}
\usepackage{hyperref}
\usepackage{enumitem}
\usepackage{algorithm}
\usepackage{algpseudocode}

\title{Why Inter-Community Events Differ in a Two-Layer Gillespie Simulator\\
and How to Restore SSA-Consistency (with a Chain-of-Cliques Testbed)}
\author{Technical report (mathematical analysis)}
\date{\today}

\begin{document}
\maketitle

\begin{abstract}
This report explains why a two-layer (macro--micro) stochastic simulator can produce different
\emph{inter-community} transmission events than a single-layer Gillespie SSA run on the full graph,
even when parameters appear ``matched.'' We provide formal root causes and propose fixes that
restore \emph{SSA-consistency} (i.e., exact equivalence in law to the full-network CTMC) for a
community-structured network, with a focus on a \emph{chain of cliques} testbed.

The analysis is grounded in the attached paper's design choices---macro rates of the form
$\lambda_{ij}(t)=\beta\,w_{ij}\,C_i(t)\,S_j(t)\,T$ and hazard-integral timing on a micro-grid, as well as
uniform seeding of imported infections at the micro layer \cite{paper}.
\end{abstract}

\section{Setup and notation}

Let $G=(V,E)$ be the full (non-aggregated) graph. Each node $v\in V$ is in one of three states
$X_v(t)\in\{S,I,R\}$. Edges may have weights $w_{uv}\ge 0$ (set $w_{uv}=1$ if unweighted).
The full-network infection/activation process is a continuous-time Markov chain (CTMC) with two
event types:

\begin{itemize}[leftmargin=*]
\item \textbf{Infection along an edge} $(u,v)\in E$: if $X_u=I$ and $X_v=S$, then $v$ becomes infected.
\item \textbf{Recovery at a node} $u$: if $X_u=I$, then $u$ recovers.
\end{itemize}

At state $x=(X_v)_{v\in V}$, the \emph{exact} SSA propensities are
\begin{align}
a_{\text{inf}}(x) &= \beta \sum_{(u,v)\in E} w_{uv}\,\mathbf{1}\{X_u=I,\,X_v=S\},\\
a_{\text{rec}}(x) &= \gamma \sum_{u\in V}\mathbf{1}\{X_u=I\}.
\end{align}
The total rate is $\Theta(x)=a_{\text{inf}}(x)+a_{\text{rec}}(x)$; the next-event time is exponential
$\Delta t\sim \mathrm{Exp}(\Theta(x))$, and the event channel is chosen with probability proportional to its
propensity.

\paragraph{Communities.}
Assume a partition $V=\bigsqcup_{i=1}^M V_i$ into communities (metanodes). Let
$E_{ij}=\{(u,v)\in E: u\in V_i,\, v\in V_j\}$ denote directed inter-community edges
($i\ne j$). In a \emph{chain of cliques} testbed, each $G[V_i]$ is a complete graph, and $E_{ij}$ exists
only for neighboring cliques in the chain.

\section{What the current two-layer model does (and where exactness is lost)}

The attached paper proposes a two-layer design \cite{paper}:

\begin{itemize}[leftmargin=*]
\item \textbf{Micro layer:} each community runs its own Gillespie SIR engine (exact within-community dynamics).
At the end of each micro step of length $\tau_{\text{micro}}$, it exports two scalars (paper: $C_i(t)=I_i(t)$
and $S_i(t)=S_i(t)$ counts \cite{paper}; your current implementation uses fractions
$\hat C_i=I_i/N_i$ and $\hat S_i=S_i/N_i$).
\item \textbf{Macro layer:} inter-community transmissions occur as CTMC events with time-dependent rates
$\lambda_{ij}(t)=\beta\,w_{ij}\,C_i(t)\,S_j(t)\,T$ \cite{paper}. Event times are sampled by a hazard-integral
scheme on a discrete time grid \cite{paper}.
\item \textbf{Import/seeding:} when a macro event $i\to j$ fires, the micro engine for community $j$ ``injects''
one infected node, in the paper drawn \emph{uniformly at random from the susceptible list}
\cite{paper}.
\end{itemize}

Even if parameters are tuned so that \emph{expected} flows roughly match, this construction is generally
\emph{not equivalent in law} to the full-network SSA on $G$. The mismatch is concentrated in the
inter-community infection channel for four interacting reasons (Sections \ref{sec:rootcauses}--\ref{sec:fixes}).

\section{Why inter-community events differ (formal root causes)}
\label{sec:rootcauses}

\subsection{Root cause A: Macro intensities generally do not equal the true SSA propensities}

In the \textbf{full-network SSA}, the instantaneous \emph{inter-community} infection propensity from
community $i$ to $j$ is exactly
\begin{equation}
\label{eq:true_lambda}
\lambda^{\star}_{ij}(x)=\beta\sum_{(u,v)\in E_{ij}} w_{uv}\,\mathbf{1}\{X_u=I,\,X_v=S\}.
\end{equation}
This is a \emph{sum over active cross-edges} (infectious source, susceptible target). It depends on the
infection states of \emph{specific boundary nodes} that participate in $E_{ij}$.

In the current \textbf{macro layer}, the rate is approximated by a low-dimensional product
\begin{equation}
\label{eq:approx_lambda}
\lambda_{ij}(t)=\beta\,w_{ij}\,C_i(t)\,S_j(t)\,T
\quad\text{(paper form)}\cite{paper}
\end{equation}
(or by using fractions $\hat C_i,\hat S_j$ in your current engine).
In general, no choice of scalars $C_i,S_j$ and constants $(w_{ij},T)$ can make \eqref{eq:approx_lambda}
equal to \eqref{eq:true_lambda} \emph{for all states}, because \eqref{eq:true_lambda} is sensitive to which
nodes are infected/susceptible \emph{among the endpoints of inter edges}.

\paragraph{Concrete chain-of-cliques example.}
Suppose clique $i$ connects to clique $i{+}1$ by a \emph{single} bridge edge $(b_i,b_{i+1})$.
Then
\[
\lambda^{\star}_{i,i+1}(x)=\beta\,\mathbf{1}\{X_{b_i}=I,\,X_{b_{i+1}}=S\}.
\]
But any macro rule based only on the \emph{fraction infected} $I_i/N_i$ and \emph{fraction susceptible}
$S_{i+1}/N_{i+1}$ necessarily assigns a \emph{strictly positive} rate as soon as $I_i>0$ and $S_{i+1}>0$,
even when $b_i$ is not yet infected. The full SSA rate is \emph{exactly zero} until the boundary node
becomes infected. This produces unavoidable timing and frequency differences of inter-community events.

\paragraph{Scaling issue with fractions.}
If you export $\hat C_i=I_i/N_i$ and $\hat S_j=S_j/N_j$, then the product $\hat C_i\hat S_j$ is
$\mathcal{O}(1)$, whereas the true propensity \eqref{eq:true_lambda} scales like the \emph{number (or total weight)
of active cross edges}. Unless $w_{ij}T$ is redefined to absorb factors such as $|E_{ij}|N_iN_j$, the macro
hazard will be systematically mis-scaled. Even with re-scaling, state-dependence on boundary nodes
still remains (previous paragraph).

\subsection{Root cause B: Loss of edge-conditioned target selection when importing infections}

In the full SSA, after selecting the \emph{inter-community} channel $(i,j)$, the simulator must still
select a \emph{specific edge} $(u,v)\in E_{ij}$ that fired, with probability
\begin{equation}
\label{eq:edge_choice}
\Pr\{(u,v)\text{ fires}\mid i\to j\} \propto
w_{uv}\,\mathbf{1}\{X_u=I,\,X_v=S\}.
\end{equation}
Equivalently, the target node $v\in V_j$ is chosen with probability proportional to the total incoming
weight from infectious neighbors in $V_i$:
\begin{equation}
\label{eq:target_choice}
\Pr\{v\text{ is infected}\mid i\to j\} \propto
\sum_{u\in V_i:(u,v)\in E_{ij}} w_{uv}\,\mathbf{1}\{X_u=I\}.
\end{equation}

By contrast, the two-layer model in the paper injects an infected node in $V_j$ by drawing it
\emph{uniformly from all susceptibles} \cite{paper}. This generally violates \eqref{eq:target_choice} because:
(i) only \emph{boundary} susceptibles incident to $E_{ij}$ should be eligible, and (ii) eligibility/weights depend on
which sources are currently infected. This changes not only \emph{who} is infected, but also \emph{future hazards},
because boundary infections disproportionately affect subsequent inter-community transmission.

For a chain of cliques with a single bridge node, uniform seeding is maximally wrong:
the full SSA always infects the \emph{unique} bridge node on the destination side, while uniform seeding infects an
interior node with probability $\approx 1-1/N_j$. This alone breaks SSA-equivalence even if event times were perfect.

\subsection{Root cause C: Macro hazard-integral on a micro-grid biases waiting times when rates jump at event times}

The hazard-integral scheme in the paper approximates
$H(t)=\int_{t_0}^t\Theta(s)\,ds$ on a discrete grid and interpolates when a unit-rate exponential budget is reached
\cite{paper}. This is appropriate when $\Theta(t)$ is smooth or at least well-approximated on the chosen grid.

However, in an exact CTMC on the full network, the inter-community rate $\Theta_{\text{inter}}(t)$ is \emph{piecewise
constant} and changes only at event times (internal infections/recoveries that flip boundary states, and inter events
themselves). Therefore, between events the correct waiting time is exponential with a \emph{constant} rate.

If the two-layer implementation updates $(C_i,S_i)$ only every $\tau_{\text{micro}}$, then the macro layer sees a
rate that is effectively \emph{time-averaged} over each interval. This produces two biases:
\begin{itemize}[leftmargin=*]
\item \textbf{Timing bias:} events that should occur shortly after a boundary node becomes infected can be delayed
until the next synchronization tick (rate too low), while events can also be triggered too early if the averaged
rate is too high.
\item \textbf{Ordering bias:} the competition between inter-community events and within-community events is altered,
changing which events happen first (and thus changing future rates recursively).
\end{itemize}

\subsection{Root cause D: ``Interrupted exponential'' and residual-budget handling can be wrong if the state changes are not treated as CTMC event times}

The paper discusses ``interrupted exponentials'' and provides a residual-budget method for deterministic imports
\cite{paper}. The key theoretical point is: \emph{any time the underlying CTMC state changes, all propensities must be
treated as updated immediately, and the next-event time must be sampled from the new state.}

If an implementation (macro or micro) does one of the following, it biases inter-event times:
\begin{itemize}[leftmargin=*]
\item carries over an exponential waiting time drawn under an old rate without conditioning or budget-update,
\item resets the hazard integral but fails to preserve the remaining budget $E-H$ correctly when rates jump,
\item handles imports/events only at the ends of fixed steps instead of at their actual firing times $t^\star$,
\end{itemize}
then the waiting-time distribution deviates from the SSA law even if the average rate is ``close.''

\section{Fixes that restore SSA-consistency}
\label{sec:fixes}

To match the full-network SSA \emph{exactly} (in distribution), you must ensure two properties:

\begin{enumerate}[leftmargin=*]
\item \textbf{Exact macro intensities:} For every pair $(i,j)$, the macro intensity equals the true inter-edge propensity
$\lambda^{\star}_{ij}(x)$ in \eqref{eq:true_lambda}.
\item \textbf{Exact edge-conditioned seeding:} Conditional on an $(i,j)$ event, the chosen infected destination node (and
source) follows the SSA edge-selection distribution \eqref{eq:edge_choice}--\eqref{eq:target_choice}.
\end{enumerate}

Once these hold, \textbf{timing} becomes straightforward: between CTMC state changes, all rates are constant,
so standard SSA exponential draws apply; no hazard-integral grid is needed for exactness.

\subsection{Fix 1: Define macro intensities as sums of true cross-edge propensities (exact)}
\label{sec:fix1}

Replace \eqref{eq:approx_lambda} by the exact definition \eqref{eq:true_lambda}:
\begin{equation}
\boxed{\;\lambda_{ij}(x)\;:=\;\lambda^{\star}_{ij}(x)=\beta\sum_{(u,v)\in E_{ij}} w_{uv}\,\mathbf{1}\{X_u=I,\,X_v=S\}\;}
\end{equation}
This is \emph{the} unique choice that reproduces the single-layer SSA channel probabilities for inter-community
events.

\paragraph{How to compute it cheaply.}
You do \emph{not} need all node-level information on the macro layer. You only need enough micro exports to
compute the active cross-edge sum. Two practical options:

\begin{itemize}[leftmargin=*]
\item \textbf{Boundary-state export (exact, minimal for chain-of-cliques).}
Let $B_i=\{u\in V_i:\exists v\notin V_i \text{ with } (u,v)\in E\}$ be boundary nodes of community $i$.
Export $X_u$ only for $u\in B_i$ (or a bitmask). Then $\lambda^{\star}_{ij}$ can be updated by inspecting only
edges $E_{ij}$, whose count is typically small.
\item \textbf{Per-neighbor boundary counts (exact for structured cross-edges).}
If $E_{ij}$ is complete bipartite between boundary subsets $B_{i\to j}\subseteq V_i$ and $B_{j\leftarrow i}\subseteq V_j$
with uniform weights, then
\begin{equation}
\lambda^{\star}_{ij}(x)=\beta\,\bigl|I\cap B_{i\to j}\bigr|\,\bigl|S\cap B_{j\leftarrow i}\bigr|.
\end{equation}
So it suffices to export the two integers $|I\cap B_{i\to j}|$ and $|S\cap B_{j\leftarrow i}|$ per neighbor.
In a chain-of-cliques with a single bridge node, these counts are $0$ or $1$.
\end{itemize}

\subsection{Fix 2: Change the import/seeding rule to match SSA edge-based selection (exact)}

When an inter-community event $i\to j$ fires, do \emph{not} infect a uniformly random susceptible in $V_j$.
Instead:
\begin{enumerate}[leftmargin=*]
\item Select a specific active cross edge $(u,v)\in E_{ij}$ with probability proportional to
$w_{uv}\mathbf{1}\{X_u=I, X_v=S\}$ (Eq.~\eqref{eq:edge_choice}).
\item Infect the \emph{specific} node $v$ in community $j$.
\end{enumerate}
This ensures that the conditional distribution of the destination node matches the full SSA exactly.

\paragraph{Chain-of-cliques consequence.}
If there is one bridge edge, the destination is deterministic: always infect the unique bridge node. This
single change eliminates a dominant source of mismatch.

\subsection{Fix 3: Replace micro-grid synchronization by an event-driven global coupling (exact)}
\label{sec:fix3}

To be SSA-consistent, the \emph{global} system must advance from one actual CTMC event time to the next.
A clean way to do this while retaining the micro/macro modularity is the \textbf{Next Reaction Method}
with a priority queue:

\begin{itemize}[leftmargin=*]
\item Each community $i$ maintains its own internal SSA clock and next internal event time $t^{(i)}_{\text{int}}$.
\item The macro layer maintains the next inter-community event time $t_{\text{inter}}$ based on the current
total inter rate $\Theta_{\text{inter}}=\sum_{i\ne j}\lambda^{\star}_{ij}(x)$.
\item The next global event time is $t^\star=\min\bigl(\min_i t^{(i)}_{\text{int}},\; t_{\text{inter}}\bigr)$.
Execute that event, update only affected propensities, and repeat.
\end{itemize}

\paragraph{Why this removes the need for hazard-integrals.}
In an exact CTMC, propensities are constant between event times. Therefore, inter-event waiting times
are exponential with constant rate; the hazard integral is only needed if you \emph{choose} to make rates vary
continuously between events (e.g., via deterministic time-varying parameters). In your current simulator,
the problematic ``time dependence'' largely comes from sampling/exporting on a grid, not from true continuous
parameter drift.

\subsection{Fix 4: Correct residual-budget handling when deterministic imports exist (exact, if you keep deterministic imports)}
If the macro layer schedules imports at deterministic calendar times $T_{\text{imp}}$ (as in the paper's micro-layer
discussion \cite{paper}), then the internal SSA in a community must decide whether the next internal event occurs
before the next import. Two exact approaches exist:

\begin{itemize}[leftmargin=*]
\item \textbf{Restart-at-import (exact, simplest):} draw the next internal event time under the current internal rate.
If an import happens earlier, apply the import, update the state (and thus internal rate), and \emph{draw anew}
from the updated state. This is exact because the CTMC is memoryless given the new state.
\item \textbf{Residual-budget method (exact, numerically stable):} keep an exponential budget $E\sim\mathrm{Exp}(1)$
and subtract used hazard $\Theta_{\text{int}}\delta$ when an import arrives first, as described in the paper \cite{paper}.
\end{itemize}

The key is to treat imports at their \emph{true firing times} and update propensities immediately.

\section{What extra micro-level state must be exported for exactness?}

A single scalar per community (infected fraction / susceptible fraction) cannot encode inter-edge propensities.
For exact SSA-consistency on a community-structured graph you minimally need, per community $i$:

\begin{enumerate}[leftmargin=*]
\item For each neighbor community $j$, enough information to compute
$\sum_{(u,v)\in E_{ij}} w_{uv}\mathbf{1}\{X_u=I, X_v=S\}$.
\end{enumerate}

A practical ``lowest sufficient'' export depends on inter-edge structure:

\begin{itemize}[leftmargin=*]
\item \textbf{General case:} export boundary node states for nodes incident to inter edges (bitmask), and maintain
the explicit list of inter edges $E_{ij}$ on the macro layer. Then updates are $O(\deg_{\text{inter}})$ per boundary
state flip.
\item \textbf{Chain of cliques (common construction):} each clique has a small boundary set (often 1--2 nodes). Exporting
the states of those nodes is sufficient and extremely cheap.
\item \textbf{Aggregated but exact for complete bipartite cross-links:} export per-neighbor counts
$|I\cap B_{i\to j}|$ and $|S\cap B_{i\to j}|$ (or weighted sums if weights vary but factorize).
\end{itemize}

\section{SSA-consistent two-layer algorithm (pseudocode)}

\begin{algorithm}[h]
\caption{Exact micro--macro SSA for inter-community transmission (SSA-consistent)}
\begin{algorithmic}[1]
\State \textbf{Precompute:} communities $V_i$; inter-edge lists $E_{ij}$; boundary nodes.
\State \textbf{Initialize:} micro states $X$; compute internal rates $\Theta^{(i)}_{\text{int}}$ for each community.
\State \textbf{Initialize inter rates:} $\lambda_{ij} \gets \beta\sum_{(u,v)\in E_{ij}} w_{uv}\mathbf{1}\{X_u=I,X_v=S\}$.
\State \textbf{Sample clocks:} for each $i$, sample $t^{(i)}_{\text{int}} \gets t + \mathrm{Exp}(\Theta^{(i)}_{\text{int}})$;
sample $t_{\text{inter}} \gets t + \mathrm{Exp}(\Theta_{\text{inter}})$ where $\Theta_{\text{inter}}=\sum_{i\ne j}\lambda_{ij}$.
\While{not finished}
  \State $t^\star \gets \min\bigl(\min_i t^{(i)}_{\text{int}},\; t_{\text{inter}}\bigr)$.
  \If{$t^\star = t_{\text{inter}}$}
     \State Choose pair $(i,j)$ with prob.\ $\lambda_{ij}/\Theta_{\text{inter}}$.
     \State Choose edge $(u,v)\in E_{ij}$ with prob.\ $\propto w_{uv}\mathbf{1}\{X_u=I,X_v=S\}$.
     \State Set $X_v \gets I$ (infect the \emph{specific} node).
     \State Update only affected internal rates in community $j$ and boundary-derived $\lambda_{\cdot\cdot}$.
     \State Resample $t^{(j)}_{\text{int}}$ from updated $\Theta^{(j)}_{\text{int}}$ (restart clock).
     \State Resample $t_{\text{inter}} \gets t^\star + \mathrm{Exp}(\Theta_{\text{inter}})$.
  \Else
     \State Let $i$ be the argmin. Execute the next internal SSA event in community $i$ at $t^\star$
            (infection on an intra edge or recovery).
     \State Update $X$ within $V_i$; if a boundary node changed state, update affected $\lambda_{ij}$ by inspecting
            its inter edges.
     \State Resample $t^{(i)}_{\text{int}} \gets t^\star + \mathrm{Exp}(\Theta^{(i)}_{\text{int}})$.
     \State \textbf{If} any $\lambda_{ij}$ changed, resample $t_{\text{inter}} \gets t^\star + \mathrm{Exp}(\Theta_{\text{inter}})$.
  \EndIf
  \State $t \gets t^\star$.
\EndWhile
\end{algorithmic}
\end{algorithm}

\paragraph{Why this matches the full-network SSA.}
The algorithm exactly simulates the same CTMC as the full SSA, but exploits structure:
intra-community events are handled by fast micro engines, while the macro engine only tracks inter-edge channels.
The key is that \emph{rates and selections are defined at the same granularity as the CTMC} (edges/nodes), so the
law of events is unchanged.

\section{If you insist on keeping a fixed micro step $\tau_{\text{micro}}$}

Exactness is generally impossible with fixed-step synchronization unless $\tau_{\text{micro}}\to 0$ and additional
assumptions hold (e.g., boundary states are well-approximated by community-wide fractions). However, you can
reduce mismatch (not eliminate it) by:

\begin{itemize}[leftmargin=*]
\item updating macro rates at \emph{every} internal micro event that changes a boundary node (event-driven within the step),
\item using the exact inter-edge sum \eqref{eq:true_lambda} computed from boundary states at those times,
\item sampling inter-event times with a hazard budget that accounts for \emph{all} rate jumps at their true times
(piecewise-constant hazard, no midpoint averaging).
\end{itemize}

But note: once you do this, you have effectively moved to the event-driven coupling of Section \ref{sec:fix3}.

\section{Summary of fixes vs.\ what each corrects}

\begin{center}
\begin{tabular}{@{}p{0.33\linewidth}p{0.60\linewidth}@{}}
\toprule
\textbf{Fix} & \textbf{Corrects} \\
\midrule
Exact macro intensities $\lambda_{ij}=\sum_{E_{ij}}$ & Makes inter-community rates equal true SSA propensities; removes wrong scaling from using fractions and removes dependence loss on boundary nodes. \\
Edge-based seeding (pick $(u,v)$) & Restores correct conditional distribution of infected node; crucial in chain-of-cliques (bridge node). \\
Event-driven coupling (global SSA/NRM) & Removes discretization/synchronization bias in waiting times and event ordering; eliminates need for midpoint hazard integration in this setting. \\
Correct residual-budget / restart handling & Prevents ``interrupted exponential'' bugs when deterministic imports occur; preserves correct waiting-time distribution after rate jumps. \\
\bottomrule
\end{tabular}
\end{center}

\section*{References}
\begin{thebibliography}{1}
\bibitem{paper}
Y.~Kuryliak and M.~Emmerich,
\emph{Towards complexity reduction of large-scale epidemic simulation in two-scale networks},
MoMLeT-2025. (Attached paper.)
\end{thebibliography}

\end{document}
